\section{Introduzione}

Questo documento illustra la progettazione e lo sviluppo di un sistema informatico di telemedicina dedicato alla gestione clinica dei pazienti affetti da diabete di tipo 2.\\
L'architettura del sistema si basa sull'interazione tra due attori principali: \\
I pazienti utilizzano la piattaforma per registrare autonomamente le misurazioni glicemiche giornaliere, documentare l'assunzione di farmaci e segnalare eventuali sintomi, comorbidità e terapie concomitanti.\\
Il medico diabetologo può, tramite la piattaforma, assegnare piani terapeutici ai suoi pazienti, oltre che compilare e consultare la loro cartella clinica e visualizzare l'andamento della glicemia.\\
Il sistema prevede una funzionalità di autenticazione centralizzata gestita dai responsabili del servizio.\\

Nelle successive sezioni sono presentati i diagrammi UML utilizzati per la progettazione e modellazione dei requisiti, seguiti da un'analisi delle decisioni strutturali e architetturali prese.
 
 \clearpage